\documentclass[]{jlreq}
% --- Engine guard for LuaLaTeX ---
\providecommand{\epTeXinputencoding}[1]{}
% --- end guard ---
\usepackage{luatexja}
\usepackage{amsmath,amssymb,amsthm}
\usepackage{tikz-cd}
\usepackage{geometry}
\usepackage{hyperref}
\geometry{margin=25mm}
\title{BTop類における初期位相と終期位相\thanks{LuaLaTeX 用に作成されています。}}
\author{CODEX (OpenAI) -- Leanファイル生成 \and CODEX (OpenAI) -- TeXファイル生成}
\date{2025年9月28日}

\theoremstyle{definition}
\newtheorem{definition}{定義}
\theoremstyle{remark}
\newtheorem{remark}{補足}

\begin{document}
\maketitle

\section{概要}
本稿の目的は、Bourbaki 流のトポロジー圏 $\mathbf{BTop}$ において、初期位相と終期位相がどのように構成され、その普遍性がどのように利用されるかを、学部生にも親しみやすい形で説明することである。
Lean 4／mathlib 上で実装された `BTopCategory.lean` を下敷きに、構造と随伴の姿を図示しながら解説する。

\section{Bourbaki 的トポロジーの考え方}
Bourbaki の立場では、位相は開集合族そのものとして与えられる。
具体的には、集合 $X$ の上の位相 $\tau$ は \emph{開集合族} $\mathcal{O}\subseteq\mathcal{P}(X)$ であり、空集合と全体集合を含み、有限交叉と任意和に閉じている。
Lean 実装では次のような構造体としてまとめられる：
\begin{definition}
\texttt{structure BourbakiTopology (X : Type u)} は、キャリア $\mathrm{carrier}$ に開集合族を持ち、空集合・全体集合・有限交叉・任意合併の安定性をフィールドとして格納する。
\end{definition}

同様に、連続写像は順序対 $(f, \text{continuity})$ として、写像と連続性の証明を合わせて一つの矢と見なす。

\section{圏 $\mathbf{BTop}$ と $\mathbf{Top}$ の同値}
Lean ファイルでは、対象を $\texttt{BTop}$、射を `BourbakiMorphism` として圏構造を実装し、Mathlib が提供する $\mathbf{Top}$ (ファイバーを持つ圏 `TopCat`) との間に忘却関手
\[\mathrm{forgetToTopCat} : \mathbf{BTop} \to \mathbf{Top}\]
を与えている。
逆向きには Mathlib の位相を Bourbaki 形式に戻す関手 `ofTopCat` があり、両者は圏論的な同値
\[\mathrm{equivTopCat} : \mathbf{BTop} \simeq \mathbf{Top}
\]
を形成する。以下の図式がすべてをまとめる：
\begin{center}
\begin{tikzcd}
\mathbf{BTop} \arrow[r, shift left=1.5, "\mathrm{forgetToTopCat}"] \arrow[d, "\mathrm{forget}"'] & \mathbf{Top} \arrow[l, shift left=1.5, "\mathrm{ofTopCat}"] \arrow[d, "\mathrm{forget}_{\mathbf{Top}}"] \\
\mathbf{Set} \arrow[r, equal] & \mathbf{Set}
\end{tikzcd}
\end{center}
縦方向は集合忘却関手であり、上段の同値により $\mathbf{Top}$ で知られている極限・余極限、随伴などが $\mathbf{BTop}$ に輸送できる。

\section{初期位相}
\subsection{構成}
$X$ を集合、$\{F_i\}_{i\in I}$ を $\mathbf{BTop}$ の対象、$\varphi_i : X \to F_i$ を集合写像の族とする。
最弱の位相で各 $\varphi_i$ を連続にするものが \emph{初期位相} である。
Lean では次の定義に一致する：
\[
\texttt{initialTopology F φ} = \texttt{ofTopologicalSpace}\Big(\bigwedge_{i\in I} \mathrm{TopologicalSpace.induced}(\varphi_i, \tau_{F_i})\Big).
\]
この位相を持つ $X$ を対象とみなしたものが `initial F φ` である。

\subsection{普遍性}
任意の $Z\in\mathbf{BTop}$ と写像 $f : Z \to X$ について、$f$ の連続性は各 $\varphi_i \circ f$ の連続性と同値である：
\[
 f \text{ が }\texttt{initial F φ} \text{ に関して連続}\quad\Longleftrightarrow\quad \forall i,~\varphi_i\circ f \text{ が }F_i\text{ へ連続}.
\]
Lean 命題 `initial_continuous_iff` がこれを形式的に表す。
普遍性により、族 $\varphi_i$ を通じて $Z$ から $X$ へ「持ち上げる」準同型 `initialLift` を構成でき、これは $f$ をそのまま関数として持つ。

\section{終期位相}
\subsection{構成}
$\{F_i\}_{i\in I}$ と写像 $\psi_i : F_i \to X$ の族を考える。
各 $\psi_i$ を連続にする最強の位相が \emph{終期位相} であり、Lean 定義は
\[
\texttt{finalTopology F ψ} = \texttt{ofTopologicalSpace}\Big(\bigvee_{i\in I} \mathrm{TopologicalSpace.coinduced}(\psi_i, \tau_{F_i})\Big)
\]
で与えられる。これにより `final F ψ` が得られる。

\subsection{普遍性}
任意の $Z$ と写像 $f : X \to Z$ について、$f$ が終期位相に関して連続であることは、各 $f\circ\psi_i$ が連続であることと同値：
\[
 f \text{ が }\texttt{final F ψ} \text{ に関して連続}\quad\Longleftrightarrow\quad \forall i,~f\circ\psi_i \text{ が }Z\text{ へ連続}.
\]
Lean 命題 `final_continuous_iff` が対応する。
これに基づき、`finalDesc` は $X$ から $Z$ への写像 $f$ を終期位相のもとで射として下降させる。

\section{随伴と図式}
初期位相は忘却関手 $\mathrm{forget} : \mathbf{BTop} \to \mathbf{Set}$ の左随伴（不変な集合写像から最弱位相を付与）を、終期位相は右随伴（最強位相を付与）を提供する。
Lean での定義 `disc_adj_forget` と `forget_adj_indisc` は、$\mathbf{BTop}$ が $\mathbf{Top}$ と同値である事実を用いて Mathlib の `TopCat.discrete` と `TopCat.trivial` の随伴を輸送したものだ。

\begin{center}
\begin{tikzcd}
\text{初期位相付与} \dashv \mathrm{forget} \dashv \text{終期位相付与}
\end{tikzcd}
\end{center}

\section{Lean 実装との対応表}
\begin{center}
\begin{tabular}{ll}
\hline
数学的概念 & Lean 中の識別子 \\
\hline
初期位相 & \texttt{initialTopology}, \texttt{initial} \\
終期位相 & \texttt{finalTopology}, \texttt{final} \\
普遍性(初期) & \texttt{initial\_continuous\_iff} \\
普遍性(終期) & \texttt{final\_continuous\_iff} \\
持ち上げ射 & \texttt{initialLift}, \texttt{initialLift\_apply} \\
下降射 & \texttt{finalDesc}, \texttt{finalDesc\_apply} \\
\hline
\end{tabular}
\end{center}

\section{例：積位相と商位相}
\begin{itemize}
  \item 直積空間の積位相は、射影写像の族に対する初期位相として得られる。
  \item 商位相は、商写像の族を終期位相として捉えると普遍性が即座に確認できる。
\end{itemize}
Lean ではこれらの構成が `initial`／`final` を土台に記述でき、普遍性クラスの証明が `simp` で片付く補題が付随している。

\section{学習のためのリソース}
\begin{itemize}
  \item Nicolas Bourbaki, \textit{Éléments de mathématique}, Topologie générale.
  \item Mathlib documentation: \url{https://leanprover-community.github.io/mathlib4_docs/}
\end{itemize}

\end{document}
